\subsection{一致地选择条带状端}\label{sec:9g}
像前面一样，作条带状端的万有选择，并考虑稳定 $d$ 叶树 $T$ 所对应的粘合构造。所得族\eref{9.12}带有两组条带状端，它们先验地可能不同。第一组通过粘合构造，继承自各 $\mathcal S^{|v|}$ 上条带状端的万有选择；第二组则通过分类映射\eref{9.13}，从 $\mathcal S^{d+1}$ 上的万有选择拉回。
\begin{definition*}
若对于每棵 $T$，都存在包含 $\{0\}\times\prod_v\mathcal R^{|v|}$ 的开子集 $\overline U\subset\overline{\mathcal R}$，使上述两组条带状端在 $U=\overline U\cap\mathcal R$ 上相同，就称条带状端的万有选择\emph{一致}。今后，将把这样的 $U$ 或 $\overline U$ 称为\emph{粘合参数充分小的子集}。
\end{definition*}
\begin{lemma}\label{st:9.3}
存在一致的条带状端万有选择。
\end{lemma}
\begin{proof}
证明简单但有些繁琐，因此只说明基本思路。首先值得注意的是，一致性条件具有归纳性质。设 $T$ 是至少有两个顶点的稳定 $d_0$ 叶树；单顶点树给出空的一致性条件。$T$ 的一致性，利用 $d=|v|-1<d_0$ 时的 $\{\epsilon_k^{d+1}\}$，部分地确定了 $\mathcal S^{d_0+1}$ 上的条带状端 $\{\epsilon_k^{d_0+1}\}$。其次，对单棵树 $T$，一致性条件能够满足。细察\stref{引理}{9.2}的证明，可见 $\overline\gamma^T$ 在粘合参数充分小的子集
\SourcePage{129}
$\overline U$ 上的限制是开嵌入，其像为 $\overline{\mathcal R}^{d_0+1}$ 中 $T$-层的一个邻域。特别地，$\gamma^T|_U$ 也是嵌入。因此，假设 $d<d_0$ 时的 $\{\epsilon_k^{d+1}\}$ 已固定，则对于某棵指定的 $T$，总能选取 $\{\epsilon_k^{d_0+1}\}$，使其沿 $\gamma^T|_U$ 的拉回，等于由粘合过程得到的 $\mathcal S|_U$ 上的条带状端。

论证的最后一步，是检查各棵 $T$ 所施加的条件彼此不矛盾。这同样是归纳论证，以 Deligne–Mumford–Stasheff 空间边界层的组合结构为基础。为具体起见，在最简单的非平凡情形 $d_0=4$ 中说明。假设已经选取 $\{\epsilon_k^3\}$ 和 $\{\epsilon_k^4\}$，并对稳定三叶树满足一致性条件。恰有五棵具有两个顶点的稳定四叶树 $T$；顶点度数必为 $3$ 和 $4$。它们给出五个延伸的粘合映射
\[
\overline\gamma^T:(-1,0]\times\mathcal R^3\times\mathcal R^4
=(-1,0]\times\{\text{点}\}\times(-1,1)\hookrightarrow\overline{\mathcal R}^5.
\]
$\overline{\mathcal R}^5$ 是五边形。限制到粘合参数充分小的子集 $\overline U_T$ 后，每个粘合映射都是嵌入，其像覆盖一个边界面的邻域，见\sr{fig:9.3}{图~9.3}。一致性要求：在这些邻域与内部 $\mathcal R^5$ 的交集上，$\{\epsilon_k^5\}$ 由 $\{\epsilon_k^3\}$、$\{\epsilon_k^4\}$ 确定。
\SourceFigure{figure-9-3}{0.40499956}{9.3}
需要证明，这五个条件可以同时满足。任意两个邻域 $\overline\gamma^{T_1}(\overline U_{T_1})$、$\overline\gamma^{T_2}(\overline U_{T_2})$，或者不交，或者在五边形的某个顶点附近重叠；该顶点对应一棵具有三个顶点的稳定四叶树 $T_3$。在第一种情形，一致性条件显然彼此不干扰；在
\SourcePage{130}
第二种情形，利用粘合映射 $\overline\gamma^{T_3}$，以及 $\{\epsilon_k^3\}$ 与 $\{\epsilon_k^4\}$ 之间的一致性，可以验证：当粘合参数充分小时，由 $T_1$、$T_2$ 粘合构造诱导的条带状端相同。因此，相应的 $\{\epsilon_k^5\}$ 的一致性条件能够同时满足。
\end{proof}

\subsection{族的模空间}\label{sec:9h}
固定带角正合辛流形 $M$。
\begin{definition*}
族 $\mathcal S\to\mathcal R$ 的一组 Lagrangian 标记，是各纤维上 Lagrangian 标记的局部常值族。与单个曲面的情形一样，选定标记后，自然为每个 $\zeta\in\Sigma^\pm$ 赋予一对正合 Lagrangian 子流形 $(L_{\zeta,0},L_{\zeta,1})$；它们是在任意纤维中相邻边界分支的标记。假设已经为 $\mathcal S\to\mathcal R$ 配备这些标记及一组条带状端，并为每对 $(L_{\zeta,0},L_{\zeta,1})$ 选定 Floer 数据 $(H_\zeta,J_\zeta)$。于是，这个族的扰动数据是一对 $(K,J)$，其中
\[
K\in\Omega^1_{\mathcal S/\mathcal R}(\mathcal S,\mathcal H),\qquad
J\in C^\infty(\mathcal S,\mathcal J),
\]
在任意纤维上的限制满足\secref{8e}的条件，包括\eref{8.7}。这个记号需要稍作说明：$K$ 是各纤维上一形式的光滑族，取值于 $M$ 上 Hamilton 函数的空间 $\mathcal H$。等价地，若 $T\mathcal S^v\to\mathcal S$ 表示沿纤维的切丛，可以把 $K$ 视为其对偶丛拉回到 $\mathcal S\times M$ 上后的一个截面，并在 $\mathcal S\times\partial M$ 附近及条带状端上附加条件。
\end{definition*}

固定 $\mathcal S\to\mathcal R$ 的扰动数据。对于每组 $\{y_\zeta\}_{\zeta\in\Sigma}$，其中 $y_\zeta\in\mathcal C(L_{\zeta,0},L_{\zeta,1})$，可以考虑模空间
\begin{equation}\label{eq:9.17}
\M_{\mathcal S}(\{y_\zeta\})=
\left\{\begin{array}{l|l}
(r,u)&r\in\mathcal R,\ u\in C^\infty(\mathcal S_r,M),\text{ 且对于扰动数据}\\[-2pt]
&\text{在纤维 }\mathcal S_r\text{ 上的限制，}u\text{ 满足\eref{8.9}、\eref{8.10}}
\end{array}\right\}.
\end{equation}
在某个小开子集 $U\subset\mathcal R$ 上，使用使条带状端保持不变的局部平凡化，可定义一个平凡 Banach 纤维丛 $\mathcal B_{\mathcal S|_U}\to U$，其纤维是\secref{8i}中的 $W^{1,p}$ 映射空间 $\mathcal B_{\mathcal S_r}$。还存在 Banach 向量丛 $\mathcal E_{\mathcal S|_U}\to\mathcal B_{\mathcal S|_U}$，在每个 $r\in U$ 上限制为 $L^p$ 空间所成的丛 $\mathcal E_{\mathcal S_r}$。非齐次 $\overline\partial$ 方程\eref{8.9}给出截面 $\mathcal B_{\mathcal S|_U}\to\mathcal E_{\mathcal S|_U}$，其零点集是所有 $\{y_\zeta\}$ 的 $\M_{\mathcal S|_U}(\{y_\zeta\})$ 的并。在模空间一点 $(r,u)$ 处，该截面的导数称为\emph{扩展线性化算子}：
\begin{equation}\label{eq:9.18}
D_{\mathcal S,r,u}:(T\mathcal B_{\mathcal S|_U})_{(r,u)}
=T_r\mathcal R\times(T\mathcal B_{\mathcal S_r})_u\longrightarrow(\mathcal E_{\mathcal S_r})_u.
\end{equation}
第二个分量正是\eref{8.17}中取 $S=\mathcal S_r$ 的通常线性化算子 $D_{S,u}$；第一个分量则是在保持 $u$ 固定时，\eref{8.9}对 $r$ 变化的导数。注意，扩展算子仍是 Fredholm 算子。对于 $(r,u)$，适当的正则性概念是 $D_{\mathcal S,r,u}$ 为满射，
\SourcePage{131}
而不是要求更强的 $D_{\mathcal S_r,u}$ 为满射；后者是将 $u$ 视为定义域固定为 $\mathcal S_r$ 的非齐次伪全纯映射时的正则性。和通常一样，若所有相关 $(r,u)$ 都满足正则性条件，就称模空间 $\M_{\mathcal S}(\{y_\zeta\})$ 或扰动数据 $(K,J)$ \emph{正则}。\znote{123}{这里要用族版本的满射性，而不是单个纤维上的满射性，因为 $r$ 本身是未知量：\eref{9.18}的定义域里含有因子 $T_r\mathcal R$，若只要求固定 $r$ 时的 $D_{\mathcal S_r,u}$ 满射，就漏掉了族沿底空间方向的形变。换句话说，这是把 $\overline\partial$ 所定义的截面与“底空间方向”作横截，属于横截性理论中“一个映射对一族参数横截”的标准情形；横截性只要求像空间被微分像张满，并不要求两个切空间互补。[来源：nLab "transversal maps"，https:/\allowbreak /\allowbreak ncatlab.org/\allowbreak nlab/\allowbreak show/\allowbreak transversal+\allowbreak maps（2026-11-15 访问）。]}

\begin{note}\label{st:9.4}
这里的表述避开了一个技术困难：只要考虑定义域变化的伪全纯映射，就会出现这个问题。对整个族 $\mathcal S\to\mathcal R$，没有典范定义的 Banach 流形 $\mathcal B_{\mathcal S}$，因为平凡化的变换在局部定义的空间 $\mathcal B_{\mathcal S|_U}$ 之间只诱导 Banach 同胚，而不诱导可微映射。然而，椭圆正则性保证，算子\eref{9.18}以及正则情形下 $\M_{\mathcal S}(\{y_\zeta\})$ 的流形结构，都不依赖于局部平凡化。
\end{note}\znote{123}{过渡映射只是 Banach 同胚而不必是微分同胚，因此无法用它们把各局部空间 $\mathcal B_{\mathcal S|_U}$ 的切丛粘成整个 $\mathcal B_{\mathcal S}$ 的切丛；这就是不存在典范 Banach 流形的原因。但\eref{9.18}仍是有意义的 Banach 丛映射，其 Fredholm 性与满射性在 Banach 同胚下拉回不变；于是零点集在每个局部平凡化下都有光滑结构，而这些结构在重叠处一致。实践中只需固定一个局部平凡化作计算，结论与选择无关。[来源：原书 §(9h)，即\stref{笔记}{9.4}的完整论证。]}

\subsection{一致地选择扰动数据}\label{sec:9i}
现在需要简要回顾粘合构造。假设有一棵 $d$ 叶树 $T$，并对每个顶点 $v\in\operatorname{Ve}(T)$，给定带条带状端和 Lagrangian 标记的 $|v|$ 标点圆盘 $S_v$。回顾旗标 $f=(v,e)$ 与 $S_v$ 的无穷远点之间存在双射。因此，在这里，每个旗标对应一对正合 Lagrangian 子流形。若对每条 $e\in\operatorname{Ed}^{\mathrm{int}}(T)$，两个旗标 $f^\pm(e)=(v^\mp(e),e)$ 对应同一对子流形，就称各 $S_v$ 的\emph{标记相容}。可以给 $\R^2\setminus T$ 的各连通分支赋予正合 Lagrangian 子流形作为标记，以直观表示这种情形，见\sr{fig:9.4}{图~9.4}。由各 $S_v$ 粘合而成的 $(d+1)$ 标点圆盘族 $\mathcal S$，以显然的方式继承 Lagrangian 标记。

进一步假设，各 $S_v$ 配备扰动数据 $(K_v,J_v)$，使每对端 $f^\pm(e)$ 上对应的 Floer 数据相同。于是，通过粘合构造，它们在 $\mathcal S$ 上确定扰动数据 $(K,J)$。这组数据光滑延伸到自然的部分紧化 $\overline{\mathcal S}$；也就是说，存在
\begin{equation}\label{eq:9.19}
\overline K\in\Omega^1_{\overline{\mathcal S}/\overline{\mathcal R}}(\overline{\mathcal S},\mathcal H),\qquad
\overline J\in C^\infty(\overline{\mathcal S},\mathcal J),
\end{equation}
在 $\mathcal S$ 上的限制为 $(K,J)$。由于 $\overline{\mathcal S}\to\overline{\mathcal R}$ 是浸没映射，在部分紧化上，逐纤维一形式的空间仍有意义。在退化程度最大的点 $0\in\overline{\mathcal R}$ 上，纤维为 $\mathcal S_0=\coprod_vS_v$，而 $(\overline K,\overline J)$ 正是给定数据 $(K_v,J_v)$ 的不交并。最后，如果粘合的是族 $\mathcal S_v$ 而不是单个曲面，上述所有结论仍然成立。

固定条带状端的万有选择。此外，对所考虑的辛流形 $M$ 中每一对正合 Lagrangian 子流形 $(L_0,L_1)$，固定 Floer 数据 $(H_{L_0,L_1},J_{L_0,L_1})$。
\begin{definition*}
\emph{扰动数据的万有选择}，是对每个 $d\geq2$ 及每个正合 Lagrangian 子流形的 $(d+1)$ 元组 $(L_0,\ldots,L_d)$，在带有这些 Lagrangian 标记的族 $\mathcal S^{d+1}\to\mathcal R^{d+1}$ 上选择扰动数据
\begin{equation}\label{eq:9.20}
(K_{L_0,\ldots,L_d},J_{L_0,\ldots,L_d}),
\end{equation}
\SourcePage{132}
使其在该族的条带状端上化为子流形对\eref{9.5}的 Floer 数据。
\end{definition*}
\SourceFigure{figure-9-4}{0.70407616}{9.4}
从现在开始，假设条带状端的万有选择是一致的。假设已经作了扰动数据的万有选择。取稳定 $d$ 叶树 $T$，将万有族 $\mathcal S^{|v|}$ 配给其各顶点，并赋予相容的 Lagrangian 标记。于是粘合后的族 $\mathcal S\to\mathcal R$ 带有两组扰动数据，它们先验地可能不同；更精确地，它们在条带状端上总是相同，但在曲面的其余部分可能不同。第一组通过粘合构造从各 $\mathcal S^{|v|}$ 上的扰动数据诱导而来；第二组通过分类映射 $\gamma^T$，从带有正确标记的 $\mathcal S^{d+1}$ 上的扰动数据万有选择拉回。
\begin{definition*}
若在所有上述情形，均满足以下条件，就称扰动数据的万有选择\emph{一致}：
\begin{itemize}
\item 存在粘合参数充分小的子集 $U\subset\mathcal R$，使上述两组扰动数据在曲面 $\mathcal S_r$（$r\in U$）的薄部相同。这里指的是粘合过程中产生的厚薄分解，见
\SourcePage{133}
\stref{笔记}{9.1}。
\item 设 $(K,J)$ 是 $\mathcal S$ 上的第一组扰动数据，即由粘合得到的数据；$(\overline K,\overline J)$ 是其在部分紧化 $\overline{\mathcal S}$ 上如\eref{9.19}的延伸。要求另一组数据，即从 $\mathcal S^{d+1}$ 拉回的数据，也光滑延伸到 $\overline{\mathcal S}$；并且在所有粘合参数都为零的子集 $\{0\}^{\operatorname{Ed}^{\mathrm{int}}(T)}\times\mathcal R^T\subset\overline{\mathcal R}$ 上，这个延伸与 $(\overline K,\overline J)$ 相同。
\end{itemize}

\end{definition*}
为了更好地理解一致性的含义，假设有序列 $r_1,r_2,\ldots\in\mathcal R^{d+1}$，收敛于 $\overline{\mathcal R}^{d+1}\setminus\mathcal R^{d+1}$ 中 $T$-层的一点。当 $k$ 充分大时，可以写成
\begin{equation}\label{eq:9.21}
r_k=\gamma^T(\{\rho_{e,k}\},\{r_{v,k}\}),
\end{equation}
其中 $r_{v,k}\in\mathcal R^{|v|}$，$v\in\operatorname{Ve}(T)$，而 $\rho_{e,k}\in(-1,0)$，$e\in\operatorname{Ed}^{\mathrm{int}}(T)$。当 $k\to\infty$ 时，每个序列 $r_{v,k}$ 收敛，且 $\rho_{e,k}$ 趋于零。更几何地，考虑带标点圆盘
\[
S_k=\mathcal S_{r_k}^{d+1},\qquad S_{v,k}=\mathcal S_{r_{v,k}}^{|v|}.
\]
那么\eref{9.21}说明：\allowbreak 按 $T$ 的边所指定的方式粘合各 $S_{v,k}$，在相邻圆盘间留下有限条带 $[0,-\log(-\rho_{e,k})/\pi]\times[0,1]$；这些条带在 $k\to\infty$ 时越来越长。所得曲面与 $S_k$ 同构。为 $S_k$ 配备某组 Lagrangian 标记 $(L_0,\ldots,L_d)$，为 $S_{v,k}$ 配备相应的相容标记。按定义，扰动数据的万有选择在每个带 Lagrangian 标记的稳定带标点圆盘上给出指定的扰动数据，因此也在这里的 $S_k$ 和 $S_{v,k}$ 上给出数据。

一致性表示：$k$ 充分大时，$S_k$ 上的扰动数据“收敛于” $\coprod_vS_{v,k}$ 上的数据。这个表述是有意义的，因为按定义，$S_k$ 是 $\coprod_vS_{v,k}$ 的某个子集的商。此外，在 $S_k$ 的有限条带以及条带状端上，这种收敛是平凡的：当 $k\gg0$ 时，那里的扰动数据就等于 $\coprod_vS_{v,k}$ 的对应部分上的数据。\znote{125}{一致性的用处是让“取极限”与“计数”可以交换：$p\leq1$ 维模空间的边界点必落在某个退化层上，按\eref{9.21}把该点附近的解展开，得到的正是 $T$ 各顶点上较低阶问题的解；只有两组扰动数据在薄部相同，这个展开才给出\eref{9.23}中所需的系数，$\Ai$ 结合律\eref{1.2}才在链层成立。条件要求对一切 $d$ 同时成立，而不是只对 $p\leq1$ 的情形，原因是同一个族既出现在低维模空间的边界上，也出现在其内部的高维模空间中。[来源：原书 §(9l) 的紧致性论证，尤其\eref{9.29}、\eref{9.30}；另一套处理见 §(10e)。]}
\begin{lemma}\label{st:9.5}
存在一致的扰动数据万有选择。
\end{lemma}
\begin{proof}
证明采用与\stref{引理}{9.3}同类的归纳论证。这里只说明证明中的一点。一般而言，曲面 $S=\mathcal S_r^{d_0+1}$ 可以同时位于不止一个粘合映射 $\gamma^T$ 的像中，而这些映射给出不同的厚薄分解。然而，可以为每棵 $T$ 选择粘合参数充分小的开子集 $U^T\subset(-1,0)^{\operatorname{Ed}^{\mathrm{int}}(T)}\times\mathcal R^T$，使得对每个 $r\in\mathcal R^{d_0+1}$，存在满足 $r\in\gamma^T(U^T)$ 的极大树 $T$。这里的极大性意指，其他具有该性质的树 $\overline T$，都能由压缩 $T$ 的边得到。条带状端选择的一致性意味着：$\gamma^T$ 所诱导的 $S$ 的厚薄分解是极大的；其他满足 $r\in\gamma^{\overline T}(U^{\overline T})$ 的 $\gamma^{\overline T}$ 所给出的分解，由这个分解忘掉若干有限薄条带而得。对于我们的一致性要求，极大
\SourcePage{134}
分解对扰动数据施加最严格的条件。

这对\stref{引理}{9.5}的证明策略有如下启示。假设已经为 $d<d_0$ 选好满足一致性条件的扰动数据\eref{9.20}。那么，对 $\mathcal S^{d_0+1}$ 的构造，应从最复杂、边数最多的树 $T$ 所对应的子集 $\gamma^T(U^T)$ 开始；它们对应于 $\overline{\mathcal R}^{d_0+1}$ 的顶点。随后按紧化空间边界层的维数递增继续构造。
\end{proof}

\begin{note}\label{st:9.6}
乍看起来，尤其将以上讨论与\secref{9g}的对应讨论比较时，似乎施加更强的一致性条件更自然：要求粘合后的族 $\mathcal S\to\mathcal R$ 上的两组扰动数据，在粘合参数充分小时完全相同。为看出它的缺点，取一棵有两个顶点、每个顶点度数均为 $3$ 的树 $T$。为每个顶点配给族 $\mathcal S^3\to\mathcal R^3$；实际上它只包含唯一的三标点圆盘 $S$。采用\eref{8.5}的记号，粘合所得的四标点圆盘为 $S\#_lS$。令所有标记都等于同一个正合 Lagrangian 子流形。如果万有扰动数据满足更强的一致性条件，它就在圆盘 $S\#_lS$ 上赋予两“半”完全相同的扰动数据；更精确地说，是厚部的两个连通分支上的数据相同。在这类扰动数据中，试图使非齐次伪全纯映射模空间正则时，就会出现类似多重覆盖的问题。我们对一致性的定义避开了这个困难。
\end{note}
